\begin{figure}[h]\centering\includegraphics[width=\columnwidth]{figures/coverage.pdf}
\caption{Validation coverage for the four retrieval configurations. The same normalized reference is checked after retrieval at depths one, three and five. Left: complete source text. Right: only passages returned by the search tool.}
\end{figure}
\subsection{Distinguishing two kinds of retrieval failure}
Reading whole documents after ranking gives a second diagnostic. At depth five, BM25 retrieved a document containing the reference for 19 questions, E5 for 13, Qwen for 17, and fusion for 26. All four runs used the same document collection and original question text. The differences between whole-document and returned-passage hits were 15, 6, 8 and 16 questions, respectively. These are differences between paired counts within a method. They do not mean that an arbitrary mention in a full document constitutes answer evidence.

At depth one, the fused result covered 12 questions under the whole-document measure and five under the passage measure. Moving from the first document to the first three raised those counts to 20 and seven. Returning five gave 26 and ten. This pattern supports exposing a separate reader: a promising document may deserve a more focused query or a different window before it is discarded. It also shows a retrieval limit that a reader cannot fix by itself. In 74 of 100 cases the normalized target string was absent from all five full documents after the initial unmodified question query.

For example, the validation question with ID 102 placed document 11266 first under fusion. Offline analysis found the normalized reference in its stored text but in none of the returned passages across the first five. The initial question includes historical events and asks for the name of an artist. A query that repeats every clue can select a passage about an event while omitting the requested attribute. A follow-up could search for the artist in context. Whether that decision actually recovers the answer must be tested by the live model.
